% =====================================================================
% MASTER_THESIS_REVISADA
% Compressão Seletiva de Contexto via Problema da Mochila:
% um Framework Não Invasivo para Extensão da Janela de Contexto em LLMs
%
% Autor:  Matheus Cardoso
% Programa: PESC / COPPE / UFRJ
%
% Projeto modular: cada capítulo é um arquivo em chapters/.
% Compilar com: tectonic main.tex
% =====================================================================
\documentclass[12pt,a4paper,oneside]{report}

\usepackage[T1]{fontenc}
\usepackage[utf8]{inputenc}
\usepackage[brazil]{babel}
\usepackage[a4paper,margin=3cm]{geometry}
\usepackage{setspace}
\usepackage{amsmath,amssymb}
\usepackage{graphicx}
\usepackage{booktabs}
\usepackage{multirow}
\usepackage{array}
\usepackage{caption}
\usepackage{enumitem}
\usepackage{tikz}
\usetikzlibrary{shapes.geometric,arrows.meta,positioning,fit,backgrounds}
\usepackage[hidelinks]{hyperref}

\onehalfspacing
\setlength{\parindent}{1.25em}
\setlength{\parskip}{0.3em}

% Rótulos de capítulo em português (report class)
\renewcommand{\bibname}{Referências Bibliográficas}

% Marcador de conteúdo pendente. Renderiza um bloco visível no PDF.
% Uso: \TODO{TESTES}{descrição do que ainda será preenchido}
\newcommand{\TODO}[2]{%
  \par\medskip\noindent
  \fcolorbox{red!60!black}{yellow!20}{%
    \parbox{0.94\linewidth}{{\ttfamily\bfseries <<TODO-#1>>}\normalfont\par\smallskip #2}%
  }\par\medskip}

\begin{document}

% =====================================================================
% CAPA / FOLHA DE ROSTO
% =====================================================================
\begin{titlepage}
\begin{center}
\textbf{\large COPPE/UFRJ}\\[0.3cm]
{\small Instituto Alberto Luiz Coimbra de Pós-Graduação e Pesquisa de Engenharia}\\[2.5cm]

{\Large\bfseries COMPRESSÃO SELETIVA DE CONTEXTO VIA PROBLEMA DA MOCHILA:\\[0.4cm]
UM \textit{FRAMEWORK} NÃO INVASIVO PARA EXTENSÃO DA JANELA DE CONTEXTO EM LLMs}\\[3cm]

{\large Matheus Cardoso}\\[3cm]

\begin{flushright}
\begin{minipage}{0.55\textwidth}
Dissertação de Mestrado apresentada ao Programa de
Pós-graduação em Engenharia de Sistemas e Computação,
COPPE, da Universidade Federal do Rio de Janeiro, como
parte dos requisitos necessários à obtenção do título de
Mestre em Engenharia de Sistemas e Computação.
\end{minipage}
\end{flushright}

\vspace{1.5cm}
\begin{flushright}
\begin{minipage}{0.55\textwidth}
Orientadores: Geraldo Xexeo\\
\phantom{Orientadores: }Claudia Susie
\end{minipage}
\end{flushright}

\vfill
Rio de Janeiro\\
Julho de 2026
\end{center}
\end{titlepage}

% =====================================================================
% RESUMO
% =====================================================================
\chapter*{Resumo}
\addcontentsline{toc}{chapter}{Resumo}

Modelos de Linguagem de Grande Escala (\textit{Large Language Models}, LLMs) tornaram-se componentes centrais de aplicações modernas, mas ainda esbarram em uma limitação estrutural, a janela de contexto. O mecanismo de atenção apresenta complexidade quadrática no comprimento da sequência e, mesmo quando a janela é nominalmente ampla, o modelo tende a não utilizar toda a informação de forma eficaz, com fenômenos como \textit{lost in the middle}, \textit{context rot} e colapso de atenção. Este trabalho investiga estratégias \textbf{não invasivas} de extensão de contexto, aplicáveis a modelos acessados via API, sem retreinamento nem modificação arquitetural, e propõe um \textit{framework} que formula a compressão de \textit{prompt} como um problema de alocação de orçamento. A partir de uma revisão sistemática das técnicas existentes e de uma avaliação empírica de cinco estratégias em \textit{benchmarks} de contexto longo, o trabalho formaliza a seleção de compressores como um problema da mochila de múltipla escolha. Dado um \textit{prompt} particionado, um orçamento efetivo de contexto e uma família de compressores parametrizados, o modelo escolhe, para cada partição, a opção de compressão que maximiza a informação preservada sob a restrição de orçamento. Os resultados empíricos indicam que estratégias simples de segmentação com sobreposição oferecem bom equilíbrio entre acurácia e latência, e sustentam a hipótese de que a compressão deve ser tratada como alocação orientada à tarefa, e não como redução uniforme de \textit{tokens}. A avaliação completa do \textit{framework} proposto é apresentada como plano experimental reprodutível.

% =====================================================================
% SUMÁRIO E LISTAS
% =====================================================================
\tableofcontents
\listoffigures
\listoftables

% =====================================================================
% CAPÍTULOS
% =====================================================================
\include{chapters/cap1_introducao}
\include{chapters/cap2_fundamentos}
\include{chapters/cap3_metodologia}
\include{chapters/cap4_resultados}
\include{chapters/cap5_conclusao}

% =====================================================================
% REFERÊNCIAS
% =====================================================================
\input{chapters/referencias}

\end{document}
